\section{Conclusioni}

KompozeR e' stato concepito come un'estensione evoluta dell'e-commerce Kompo, con l'obiettivo di rendere piu' semplice e guidata la configurazione e l'acquisto di scaffalature personalizzabili. Il progetto affronta il problema sia dal punto di vista dell'esperienza utente, introducendo un configuratore visuale, un chatbot di supporto e notifiche contestuali, sia dal punto di vista dell'organizzazione software, tramite una struttura modulare basata su servizi distinti.

\subsection{Risultati raggiunti}

In questa sottosezione conviene riassumere i principali risultati del lavoro svolto, ad esempio:

\begin{itemize}
	\item definizione dei requisiti e del dominio applicativo;
	\item progettazione di un'architettura web coerente con il problema affrontato;
	\item individuazione di servizi separati per autenticazione, catalogo, configurazione, carrello, notifiche, chatbot e reporting;
	\item definizione di contratti applicativi chiari per API, messaggi real-time e DTO di dominio.
\end{itemize}

\subsection{Limiti attuali}

E' utile esplicitare anche i limiti del progetto nella sua fase corrente, ad esempio eventuali funzionalita' solo progettate, copertura di test ancora parziale o assenza di alcune ottimizzazioni implementative. Una sezione di questo tipo rende il report piu' rigoroso e chiarisce bene il perimetro del lavoro effettivamente svolto.

\subsection{Sviluppi futuri}

Tra le evoluzioni naturali del progetto rientrano il consolidamento dell'implementazione, il completamento delle attivita' di verifica e il raffinamento delle funzionalita' piu' avanzate lato utente e lato amministrativo. In questa parte puoi anche richiamare il fatto che gli aspetti piu' strettamente distribuiti saranno approfonditi separatamente nel report DS.

\subsection{Valutazione finale}

Nel complesso, il progetto rappresenta un caso applicativo interessante per il contesto ASW perche' combina progettazione di servizi web, organizzazione modulare del backend, attenzione all'interazione utente e definizione accurata dei contratti tra componenti. Il report mette quindi in evidenza soprattutto la qualita' della soluzione come applicazione e servizio web, lasciando a un documento separato l'analisi dei temi distribuiti piu' avanzati.